% Slides, 16:9. Build with pdfLaTeX; add the class option "handout" for a printable version without overlays.
\documentclass[aspectratio=169,11pt]{beamer}

\usepackage[T1]{fontenc}
\usepackage[utf8]{inputenc}
\usepackage[sfdefault]{FiraSans}
\usepackage{newtxsf}
\usepackage{booktabs}
\usepackage{tikz}
\usepackage{appendixnumberbeamer}

\usetheme[progressbar=frametitle,numbering=fraction,block=fill]{metropolis}
\setbeamertemplate{bibliography item}[text]

\title{Title of the Talk}
\subtitle{Event, City, Date}
\author{Norbert Micheľ}
\institute{Institute of Computer Science, Pavol Jozef Šafárik University in Košice}
\date{}

\begin{document}

\maketitle

\begin{frame}{Outline}
  \tableofcontents
\end{frame}

\section{Statements}

\begin{frame}{Definition and theorem}
  \begin{definition}
    The \alert{degree} $\deg v$ of a vertex $v$ of a graph $G=(V,E)$ is the number of edges incident with $v$.
  \end{definition}
  \pause
  \begin{theorem}
    Every graph $G=(V,E)$ satisfies $\sum_{v\in V}\deg v = 2|E|$.
  \end{theorem}
  \pause
  \begin{proof}
    Every edge is counted once at each of its two ends.
  \end{proof}
  \note{Speaker notes go here; they appear with \texttt{\textbackslash setbeameroption\{show notes\}}.}
\end{frame}

\section{Figures and tables}

\begin{frame}{Figure built step by step}
  \centering
  \begin{tikzpicture}[v/.style={circle,fill=black,inner sep=1.6pt},every edge/.style={draw,thick}]
    \path[use as bounding box] (-0.5,-1.6) rectangle (4.5,1.6);
    \node[v,label=left:$s$] (s) at (0,0) {};
    \node[v] (a) at (2,1) {};
    \node[v] (b) at (2,-1) {};
    \node[v,label=right:$t$] (t) at (4,0) {};
    \path (s) edge (a) (s) edge (b) (a) edge (b) (a) edge (t) (b) edge (t);
    \draw<2->[mLightBrown,ultra thick] (s) -- (a) -- (t);
    \draw<3->[mDarkTeal,dashed,thick] (1,1.5) -- (1,-1.5);
  \end{tikzpicture}

  \medskip
  \begin{overlayarea}{\textwidth}{1.5em}
    \centering
    \only<2>{An $s$--$t$ path.}\only<3>{The edges that cross the dashed line form an $s$--$t$ cut.}
  \end{overlayarea}
\end{frame}

\begin{frame}{Table}
  \centering
  \begin{tabular}{lrr}
    \toprule
    Instance & $n$ & Value\\
    \midrule
    A & 10 & 1\\
    B & 20 & 2\\
    \bottomrule
  \end{tabular}
\end{frame}

\begin{frame}{Two columns}
  \begin{columns}[T]
    \begin{column}{0.48\textwidth}
      \begin{itemize}
        \item<1-> First point
        \item<2-> Second point
      \end{itemize}
    \end{column}
    \begin{column}{0.48\textwidth}
      \begin{block}{Remark}
        Blocks and lists share a frame in columns.
      \end{block}
    \end{column}
  \end{columns}
\end{frame}

\begin{frame}{Summary}
  \begin{itemize}
    \item Main result in one line.
    \item Open question in one line.
  \end{itemize}
\end{frame}

\begin{frame}[standout]
  Thank you
\end{frame}

\appendix

\begin{frame}{References}
  \nocite{NagamochiKameda1996}
  \bibliographystyle{plain}
  \bibliography{references}
\end{frame}

\end{document}
